%REVIEW-ADDR: W1-zvf-tautology
%REVIEW-ADDR: Q1-zvf-formalization
\section{Formal Definition of the Zero-Variance Fraction, Non-Tautology, and Scope}
\label{sec:appendix:zvf-formalization}

\subsection{Formal Definition}

Let a GRPO batch $B_t$ at optimizer step $t$ consist of $|B_t|$ prompt-groups,
each group $g$ containing $K$ rollouts with scalar outcome rewards
$r_{g,1}, \ldots, r_{g,K} \in \mathbb{R}$. We define the \emph{Zero-Variance
Fraction} (ZVF) as
\begin{equation}
\label{eq:zvf}
\mathrm{ZVF}_t \;=\; \frac{1}{|B_t|} \sum_{g \in B_t}
  \mathbb{1}\!\left[\;\widehat{\mathrm{Var}}(r_{g,1}, \ldots, r_{g,K}) \le \varepsilon\;\right],
\end{equation}
where $\widehat{\mathrm{Var}}$ is the unbiased sample variance and
$\varepsilon$ is a small numerical tolerance
(we use $\varepsilon = 10^{-6}$ for rewards normalized to $[0,1]$).
Intuitively, $\mathrm{ZVF}_t$ is the fraction of groups at step $t$ whose
rollouts all collapsed to the same outcome; those groups contribute zero
group-relative advantage and thus zero gradient signal to the GRPO update.

\subsection{Why ZVF Is Not a Tautological Re-expression of Mean Reward}

A reviewer may reasonably suspect that under binary outcome rewards
$r \in \{0,1\}$ and group-relative advantages, $\mathrm{ZVF}_t$ is a direct
function of the batch mean reward $\bar{r}_t$: if every group succeeds
(or fails) uniformly, $\mathrm{ZVF}_t = 1$. This is true at the extremes
$\bar{r}_t \in \{0,1\}$ but does \emph{not} hold at intermediate values.
For $K$ i.i.d.~Bernoulli rollouts with prompt-specific success probability
$p_g$, the expected ZVF under a null independence model is
\begin{equation}
\mathbb{E}[\mathrm{ZVF}_t] \;=\; \mathbb{E}_{p_g}\!\left[\,p_g^{K} + (1-p_g)^{K}\,\right].
\end{equation}
The right-hand side depends on the full prompt-difficulty distribution
$\{p_g\}$, not only on its first moment. Two populations can therefore
share the same mean reward $\bar r = 0.5$ yet have very different ZVF:
a bimodal population with $p_g \in \{0,1\}$ yields $\mathrm{ZVF}=1$,
while a uniform-hard population with all $p_g = 0.5$ and $K=8$ yields
$\mathrm{ZVF} \approx 2 \cdot (0.5)^8 \approx 0.008$. The substantive
claim is thus not ``high mean reward implies high ZVF,'' but rather that
ZVF is sensitive to \emph{how} success mass is distributed across prompts.

\subsection{What the Released Artifact Establishes}

The anonymous artifact bundled with this paper contains aggregate reward
trajectories, held-out checkpoint summaries, and the cross-framework parser
specification in Appendix~\ref{sec:appendix:zvf-pipeline}. It does
\emph{not} contain the per-group step-level rollout logs needed to support
a new partial-correlation table that controls simultaneously for
batch-mean reward, entropy, KL drift, and advantage variance. We therefore
do \emph{not} claim a fresh incremental-$R^2$ estimate here, and we
withdraw any implication that the current artifact proves a new
partial-correlation result beyond the already released summary traces.

The narrower empirical claim made in the revised paper is the one directly
supported by the released evidence: in outcome-reward GRPO, high ZVF
coincides with the disappearance of within-group contrast, and those are
exactly the runs where reward traces plateau or regress because the
group-relative advantage has collapsed to zero. This is the operational
meaning of ZVF in the present benchmark.

\subsection{Boundary Conditions and Scope}

We state explicitly when ZVF is \emph{not} informative, partly in response
to the tautology concern:
\begin{itemize}[leftmargin=1.2em,itemsep=0pt]
  \item \textbf{Dense or continuous rewards.} When rewards are continuous
        (e.g.,~process-reward models or graded code-execution partial credit),
        $\widehat{\mathrm{Var}}$ is almost surely nonzero and
        $\mathrm{ZVF} \to 0$. ZVF degenerates and must be replaced by a
        per-step-reward-variance diagnostic or the surrogates discussed in
        Section~\ref{sec:extended-related-work}.
  \item \textbf{$K=1$.} Group variance is undefined and ZVF is trivially $1$.
  \item \textbf{SFT-saturated baselines.} When the policy already solves
        the task ($\bar p_g \to 1$ for essentially every prompt),
        $\mathrm{ZVF} \to 1$ and the diagnostic loses discriminative power;
        in this regime RL no longer provides useful signal anyway.
  \item \textbf{Non-outcome-reward RL.} For DPO- or regression-style
        training there are no rollout groups and ZVF is ill-defined.
\end{itemize}
Within its intended scope -- outcome-reward GRPO on verifiable tasks with
$K \ge 4$ and non-saturated reward support -- ZVF is a cheap, model-agnostic
early-warning diagnostic. Outside that scope, the paper now recommends
explicit surrogates instead of over-claiming portability.
